\documentclass[twocolumn,10pt]{article}

\usepackage[margin=0.75in]{geometry}
\usepackage{amsmath,amssymb,amsthm}
\usepackage{mathtools}
\usepackage{bm}
\usepackage{graphicx}
\usepackage{booktabs}
\usepackage{hyperref}
\usepackage{cleveref}
\usepackage{float}
\usepackage{caption}
\usepackage{subcaption}
\usepackage{xcolor}
\usepackage{algorithm}
\usepackage{algorithmic}
\usepackage{enumitem}
\usepackage{physics}
\usepackage{tikz}
\usetikzlibrary{shapes,arrows,positioning,calc,decorations.pathmorphing}
\usepackage{pgfplots}
\pgfplotsset{compat=1.18}
\usepackage{siunitx}
\usepackage[numbers,sort&compress]{natbib}

% Theorem environments
\newtheorem{theorem}{Theorem}[section]
\newtheorem{lemma}[theorem]{Lemma}
\newtheorem{proposition}[theorem]{Proposition}
\newtheorem{corollary}[theorem]{Corollary}
\newtheorem{definition}[theorem]{Definition}
\newtheorem{axiom}{Axiom}
\theoremstyle{remark}
\newtheorem{remark}[theorem]{Remark}
\newtheorem{example}[theorem]{Example}
\newtheorem{principle}[theorem]{Principle}

% Custom commands
\newcommand{\kB}{k_{\mathrm{B}}}
\newcommand{\Depth}{\mathcal{M}}
\newcommand{\Partition}{\mathcal{P}}
\newcommand{\Sk}{S_k}
\newcommand{\St}{S_t}
\newcommand{\Se}{S_e}
\newcommand{\eps}{\varepsilon}
\newcommand{\Cost}{\mathcal{E}}
\newcommand{\Mfree}{M_{\mathrm{free}}}
\newcommand{\Mbound}{M_{\mathrm{bound}}}
\newcommand{\Mmax}{M_{\mathrm{max}}}
\newcommand{\Mmin}{M_{\mathrm{min}}}
\newcommand{\Gres}{\mathcal{G}}
\newcommand{\Scoord}{\mathbf{S}}

\hypersetup{
    colorlinks=true,
    linkcolor=blue,
    citecolor=blue,
    urlcolor=blue,
    pdftitle={Membrane-Mediated Categorical Conversion},
    pdfauthor={Kundai Farai Sachikonye}
}

\title{\textbf{Membrane-Mediated Categorical Conversion: \\Non-Quantizable Processors via Partition Extinction on Binary Substrates}}

\author{
Kundai Farai Sachikonye\\
AIMe Registry for Artificial Intelligence\\
\texttt{kundai.sachikonye@bitspark.com}
}

\date{}

\begin{document}

\maketitle

\begin{abstract}
\noindent We introduce the categorical converter---a processor architecture that uses virtual membranes to convert binary digital computation into continuous categorical computation without intermediate analog stages. The architecture rests on three pillars derived from first principles. First, we prove that any binary processor is a state counter operating at base $b = 2$, generating entropy $\Delta S = \kB \ln 2$ per clock transition, and that this state counting is isomorphic to partition traversal in bounded phase space. Second, we establish the Partition Extinction Theorem: when temperature drops below a critical threshold $T_c$ for a group of binary states, those states become categorically indistinguishable and merge into a single continuous categorical state---the binary quantization dissolves. Third, we prove that virtual membranes---boundary conditions in partition space implemented through oscillatory dynamics---serve as transducers between the binary and categorical regimes, with the capacity formula $C(n) = 2n^2$ providing the natural bijection between binary state counts and ternary partition depth. The resulting processor has no fixed quantizable unit: its computational resolution is a continuous function of temperature, ranging from fully discrete (high $T$, binary) through ternary ($T \approx T_c$) to fully continuous (low $T$, categorical). We derive the complete thermodynamics of categorical conversion, prove that the conversion is reversible with efficiency bounded by Carnot, and establish that a membrane of area $A$ implements a categorical processor with $\sim 10^{20}$ partition operations per second---exceeding any binary supercomputer. The framework unifies digital and analog computation as limiting cases of a single partition-depth continuum, resolves the discretization paradox in computational physics, and provides a physical foundation for biological computation. We present experimental predictions including temperature-dependent computational resolution, membrane-mediated error correction, and the categorical speedup: problems requiring $O(n)$ binary steps require only $O(\log n)$ categorical navigations.
\end{abstract}

%==============================================================================
\section{Introduction}
\label{sec:introduction}
%==============================================================================

\subsection{The Quantization Assumption}

All conventional processors rest on a single assumption: computation requires quantization. The bit---a system in one of exactly two distinguishable states---is the atomic unit of digital computation \cite{shannon1948}. Logic gates combine bits through Boolean operations. Clock cycles synchronize state transitions. The entire edifice of computer science, from Turing machines \cite{turing1936} to modern GPUs, is built on the premise that information must be discretized into countable, combinable units before it can be processed.

This assumption has been spectacularly productive. Moore's law has delivered exponential growth in binary processing power for six decades \cite{moore1965}. Yet the assumption is not a law of nature---it is an engineering choice. Nature computes without bits. Protein folding, neural signaling, gene regulation, and cellular decision-making all proceed through continuous dynamics, not discrete logic \cite{bray1995,hartwell1999}. The question is whether these biological computations are merely ``messy implementations'' of underlying digital logic, or whether they represent a fundamentally different---and potentially more powerful---computational paradigm.

We prove the latter. We show that binary digital computation and continuous categorical computation are limiting cases of a single framework governed by partition depth, and that membranes---physical or virtual---serve as the natural transducers between these regimes.

\subsection{The Bounded Phase Space Framework}

We work within the Bounded Phase Space (BPS) framework, which begins from a single axiom:

\begin{axiom}[Bounded Phase Space Law]
\label{axiom:bps}
All physical systems occupy bounded regions of phase space admitting partition and nesting.
\end{axiom}

From this axiom, the following results have been established:

\begin{enumerate}[label=(\roman*)]
    \item \textbf{Partition coordinates} $(n, \ell, m, s)$: hierarchical state labeling with principal depth $n$, angular complexity $\ell$, orientation $m$, chirality $s$.
    \item \textbf{Capacity formula}: $C(n) = 2n^2$ distinguishable states at depth $n$.
    \item \textbf{Partition depth}: $\Depth = \sum_i \log_b(k_i)$, the hierarchical structure required for state distinction, where $b$ is the encoding base.
    \item \textbf{Composition Theorem}: Binding reduces partition depth; the deficit is released as energy.
    \item \textbf{Compression Theorem}: The cost of distinguishing $N$ states in volume $\mathcal{V}$ diverges as $\Cost \propto N \ln(N\mathcal{V}_0/\mathcal{V})$.
    \item \textbf{Partition Extinction Theorem}: When partition operations between entities become undefined, the associated transport coefficient vanishes exactly.
    \item \textbf{State counting}: Temporal evolution IS state counting---$dM/dt = 1/\langle\tau_p\rangle$, where $M$ is the partition count and $\tau_p$ is the partition duration.
    \item \textbf{S-entropy coordinates}: Three information-theoretic dimensions $(\Sk, \St, \Se) \in [0,1]^3$ encoding knowledge, temporal, and evolutionary entropy.
\end{enumerate}

\subsection{Summary of Main Results}

\begin{enumerate}
    \item \textbf{Binary Processor as State Counter} (\S\ref{sec:binary}): Every binary processor is a base-2 partition traversal system. Clock cycles are partition transitions. Bits are partition coordinates at depth 1.

    \item \textbf{The Categorical Conversion Theorem} (\S\ref{sec:conversion}): Temperature-mediated partition extinction converts binary states into continuous categorical states. The conversion is governed by the partition base transformation $b: 2 \to 3 \to \infty$.

    \item \textbf{The Membrane Transducer} (\S\ref{sec:membrane}): Virtual membranes implement the boundary conditions for categorical conversion. The capacity formula $C(n) = 2n^2$ provides the bijection between binary counts and ternary depth.

    \item \textbf{The Non-Quantizable Processor} (\S\ref{sec:processor}): The categorical processor has no fixed computational unit. Its resolution is a continuous function of thermodynamic state.

    \item \textbf{Categorical Thermodynamics} (\S\ref{sec:thermodynamics}): The complete thermodynamics of categorical conversion, including efficiency bounds, entropy production, and reversibility conditions.

    \item \textbf{The Categorical Speedup} (\S\ref{sec:speedup}): Navigation in S-entropy space achieves $O(\log n)$ complexity for problems requiring $O(n)$ binary steps.

    \item \textbf{Biological Implementation} (\S\ref{sec:biological}): Cell membranes are natural categorical converters, explaining the computational power of biological systems.
\end{enumerate}

%==============================================================================
\section{The Binary Processor as State Counter}
\label{sec:binary}
%==============================================================================

\subsection{Binary Computation in the Partition Framework}

\begin{theorem}[Binary Processor Identity]
\label{thm:binary_identity}
Any binary digital processor with clock frequency $f$ and word length $w$ is a state counter operating at base $b = 2$ with:
\begin{enumerate}
    \item Partition depth per clock cycle: $\Depth_{\text{cycle}} = w$
    \item Entropy production per cycle: $\Delta S_{\text{cycle}} = w \kB \ln 2$
    \item State counting rate: $R = f \cdot w$ partition transitions per second
    \item Total partition depth traversed in time $t$: $\Depth(t) = f \cdot w \cdot t$
\end{enumerate}
\end{theorem}

\begin{proof}
\textbf{Part 1: Bits as partition coordinates.}
A single bit is a system with exactly two distinguishable states: $\{0, 1\}$. In the partition framework, this is a partition at depth $n = 1$ with capacity $C(1) = 2 \times 1^2 = 2$. The partition coordinates are $(n=1, \ell=0, m=0, s=\pm 1/2)$, where the chirality $s$ encodes the bit value:
\begin{equation}
    \text{bit} = \begin{cases} 0 & \text{if } s = -1/2 \\ 1 & \text{if } s = +1/2 \end{cases}
\end{equation}

The partition depth of a single bit is:
\begin{equation}
    \Depth_{\text{bit}} = \log_2(2) = 1
\end{equation}

\textbf{Part 2: Word as partition state.}
A $w$-bit word occupies a partition state with $2^w$ possible values. The partition depth is:
\begin{equation}
    \Depth_{\text{word}} = \log_2(2^w) = w
\end{equation}

This corresponds to a compound partition coordinate $(n, \ell, m, s)^{\otimes w}$---the tensor product of $w$ single-bit partitions.

\textbf{Part 3: Clock cycle as partition transition.}
Each clock tick potentially changes the processor state from one $w$-bit configuration to another. By the state counting framework, this is a partition transition---the system moves from one cell in partition space to an adjacent cell.

The entropy produced per transition follows from the depth-entropy equivalence:
\begin{equation}
    \Delta S = \kB \ln b \cdot \Delta\Depth = \kB \ln 2 \cdot w
\end{equation}

For a 64-bit processor at $f = 3$ GHz:
\begin{equation}
    R = 3 \times 10^9 \times 64 = 1.92 \times 10^{11} \text{ partition transitions/s}
\end{equation}
\begin{equation}
    \dot{S} = 1.92 \times 10^{11} \times 1.38 \times 10^{-23} \times \ln 2 \approx 1.84 \times 10^{-12} \text{ W/K}
\end{equation}

\textbf{Part 4: Cumulative depth.}
Over time $t$, the total partition depth traversed is:
\begin{equation}
    \Depth(t) = R \cdot t = f \cdot w \cdot t
\end{equation}

In one second: $\Depth(1) = 1.92 \times 10^{11}$. This is the processor's ``partition odometer''---how far it has traveled through partition space.
\end{proof}

\subsection{The Quantization Constraint}

\begin{theorem}[Binary Quantization Constraint]
\label{thm:quantization}
A binary processor is constrained to partition depth increments of exactly $\Delta\Depth = 1$ (one bit) per elementary operation. It cannot access partition depths between integers.
\end{theorem}

\begin{proof}
Each bit transition changes exactly one partition coordinate from $s = -1/2$ to $s = +1/2$ or vice versa. The depth change is:
\begin{equation}
    \Delta\Depth_{\text{bit}} = |\log_2(2) - \log_2(1)| = 1
\end{equation}

To access a partition depth of, say, $\Depth = 1.585$ (which corresponds to $\log_2 3$---a ternary state), the binary processor must use \emph{multiple} bits:
\begin{equation}
    \lceil \log_2 3 \rceil = 2 \text{ bits to encode 3 states}
\end{equation}

But 2 bits encode 4 states, not 3. The processor wastes $2^2 - 3 = 1$ state---25\% overhead. This is the \emph{quantization waste}: binary encoding of non-power-of-2 partition structures necessarily wastes states.

More generally, encoding $N$ states in binary requires $\lceil \log_2 N \rceil$ bits, wasting $2^{\lceil \log_2 N \rceil} - N$ states. The waste fraction:
\begin{equation}
    W(N) = 1 - \frac{N}{2^{\lceil \log_2 N \rceil}}
\end{equation}

For ternary ($N = 3$): $W = 25\%$. For 5-state ($N = 5$): $W = 37.5\%$. For 7-state ($N = 7$): $W = 12.5\%$.

The average waste across all $N$ is:
\begin{equation}
    \langle W \rangle = 1 - \frac{1}{\ln 2} \approx 0.31
\end{equation}

A binary processor wastes approximately 31\% of its partition capacity on average when encoding non-binary partition structures.
\end{proof}

\subsection{The Landauer Limit and Binary Entropy}

\begin{proposition}[Binary Landauer Bound]
\label{prop:landauer}
The minimum energy dissipated per binary partition transition at temperature $T$ is:
\begin{equation}
    E_{\min} = \kB T \ln 2
\end{equation}
This is the Landauer limit \cite{landauer1961}, here derived as a consequence of partition depth change.
\end{proposition}

\begin{proof}
A binary transition changes partition depth by $\Delta\Depth = 1$ at base $b = 2$. By the depth-entropy equivalence:
\begin{equation}
    \Delta S = \kB \ln 2 \cdot \Delta\Depth = \kB \ln 2
\end{equation}

The minimum energy dissipated is:
\begin{equation}
    E_{\min} = T \Delta S = \kB T \ln 2
\end{equation}

At $T = 300$ K: $E_{\min} = 2.87 \times 10^{-21}$ J $= 0.018$ eV.

Current processors dissipate $\sim 10^{-15}$ J per transition---$10^6$ times the Landauer limit. The gap is due to non-ideal switching, leakage, and interconnect losses \cite{frank2005}.
\end{proof}

\begin{corollary}[Base-Dependent Landauer Limit]
\label{cor:base_landauer}
For a processor operating at base $b$, the Landauer limit is:
\begin{equation}
    E_{\min}(b) = \kB T \ln b
\end{equation}
A ternary transition ($b = 3$) dissipates $\kB T \ln 3 \approx 1.585 \kB T \ln 2$---58.5\% more energy per transition but encodes $\log_2 3 \approx 1.585$ bits of information. The information per unit energy is the same:
\begin{equation}
    \frac{I}{E} = \frac{\log_2 b}{\kB T \ln b} = \frac{1}{\kB T \ln 2} = \text{const}
\end{equation}

The Landauer efficiency is base-independent. There is no thermodynamic advantage to any particular base---the advantage of ternary and continuous encoding lies in \emph{structural matching}, not energy efficiency.
\end{corollary}

\begin{figure*}[!htbp]
    \centering
    \includegraphics[width=0.9\textwidth]{panels/panel_exp09_binary_entropy.png}
    \caption{\textbf{Binary quantisation waste in the capacity formula $C(n) = 2n^2$.}
    (\textbf{A}) Percentage of wasted binary capacity for partition levels $n = 1$ to $30$. Since $C(n) = 2n^2$ is not a power of two for most $n$, encoding in $\lceil\log_2 C(n)\rceil$ bits wastes a fraction $1 - C(n)/2^{\lceil\log_2 C(n)\rceil}$ of the address space. The horizontal dashed line marks the capacity-weighted mean waste of $\sim 29\%$, approaching the theoretical $\sim 31\%$ limit.
    (\textbf{B}) Encoding efficiency for binary (blue) and ternary (red) representations versus partition level. Both oscillate between $0.5$ and $1.0$ as $C(n)$ periodically aligns with or misaligns from integer powers of the base, with neither system consistently dominating---motivating the use of a continuous effective base.
    (\textbf{C}) 3D scatter of partition level $n$, binary partition depth $M_2 = \log_2 C(n)$, and ternary partition depth $M_3 = \log_3 C(n)$, coloured by binary waste percentage. The strict linear relationship $M_3 = M_2 \cdot \log_2 3^{-1}$ confirms the universal conversion ratio, while the colour variation reveals waste as a purely discretisation artefact.
    (\textbf{D}) Step plot comparing the natural capacity $C(n)$ (blue) with the binary ceiling $2^{\lceil\log_2 C(n)\rceil}$ (red dashed). The vertical gaps between the curves directly visualise the wasted states at each level, demonstrating that binary encoding systematically over-allocates address space for non-power-of-two state counts.}
    \label{fig:exp09_binary_waste}
\end{figure*}

%==============================================================================
\section{The Categorical Conversion Theorem}
\label{sec:conversion}
%==============================================================================

\subsection{Temperature and Partition Distinguishability}

\begin{definition}[Partition Temperature]
\label{def:partition_temp}
The \emph{partition temperature} $T_\Partition$ of a group of states is the temperature at which thermal fluctuations equal the energy difference between states:
\begin{equation}
    \kB T_\Partition = \Delta E_{\text{states}}
\end{equation}
For $T > T_\Partition$: states are thermally distinguishable (partition operations are defined).
For $T < T_\Partition$: states are thermally indistinguishable (partition operations become undefined).
\end{definition}

\begin{theorem}[Temperature-Controlled Distinguishability]
\label{thm:temp_distinguish}
The effective number of distinguishable states $N_{\text{eff}}$ in a group of $N$ states with energy spacing $\Delta E$ at temperature $T$ is:
\begin{equation}
    N_{\text{eff}}(T) = 1 + (N-1) \cdot \left[1 - \exp\left(-\frac{\Delta E}{\kB T}\right)\right]
\end{equation}
\end{theorem}

\begin{proof}
Consider $N$ states with energies $E_0, E_0 + \Delta E, E_0 + 2\Delta E, \ldots, E_0 + (N-1)\Delta E$.

At temperature $T$, the Boltzmann probability of state $k$ is:
\begin{equation}
    p_k = \frac{e^{-k\Delta E / \kB T}}{\sum_{j=0}^{N-1} e^{-j\Delta E / \kB T}}
\end{equation}

Two states $k$ and $k'$ are distinguishable if $|p_k - p_{k'}| > \delta$ for some resolution threshold $\delta$. For adjacent states:
\begin{equation}
    \frac{p_k}{p_{k+1}} = e^{\Delta E / \kB T}
\end{equation}

When $\Delta E \gg \kB T$: ratio $\to \infty$, states perfectly distinguishable, $N_{\text{eff}} = N$.
When $\Delta E \ll \kB T$: ratio $\to 1$, states indistinguishable, $N_{\text{eff}} = 1$.

The interpolation:
\begin{equation}
    N_{\text{eff}} = 1 + (N-1)(1 - e^{-\Delta E/\kB T})
\end{equation}
satisfies both limits and gives the effective number of distinguishable partition states.
\end{proof}

\subsection{The Effective Partition Base}

\begin{definition}[Effective Base]
\label{def:eff_base}
The effective partition base at temperature $T$ is:
\begin{equation}
    b_{\text{eff}}(T) = 1 + (b_{\max} - 1)(1 - e^{-\Delta E / \kB T})
\end{equation}
where $b_{\max}$ is the maximum base supported by the physical substrate.
\end{definition}

This is the central quantity of the paper. It defines a \emph{continuous} variable that interpolates between:
\begin{itemize}
    \item $b_{\text{eff}} = 1$: No partition operations possible (partition extinction)
    \item $b_{\text{eff}} = 2$: Binary (standard digital computation)
    \item $b_{\text{eff}} = 3$: Ternary (the natural base of the S-entropy coordinate system)
    \item $b_{\text{eff}} \to \infty$: Continuous (fully categorical computation)
\end{itemize}

\begin{theorem}[Categorical Conversion Theorem]
\label{thm:categorical_conversion}
A binary processor with state energy spacing $\Delta E$ can be converted to a categorical processor operating at effective base $b_{\text{eff}}$ by controlling the temperature $T$ according to:
\begin{equation}
    T = -\frac{\Delta E}{\kB \ln\left(1 - \frac{b_{\text{eff}} - 1}{b_{\max} - 1}\right)}
\end{equation}

The conversion passes through three regimes:
\begin{enumerate}
    \item \textbf{Binary regime} ($T \gg \Delta E / \kB$): $b_{\text{eff}} \approx b_{\max}$. All states distinguishable. Standard digital computation.
    \item \textbf{Ternary regime} ($T \approx T_3$): $b_{\text{eff}} = 3$. The natural base for S-entropy navigation. Optimal categorical computation.
    \item \textbf{Extinction regime} ($T \ll \Delta E / \kB$): $b_{\text{eff}} \to 1$. States merge. Partition operations become undefined. The computational substrate becomes a single categorical entity.
\end{enumerate}

The ternary crossover temperature is:
\begin{equation}
    T_3 = -\frac{\Delta E}{\kB \ln\left(1 - \frac{2}{b_{\max} - 1}\right)}
\end{equation}
\end{theorem}

\begin{proof}
The proof follows directly from Definition~\ref{def:eff_base} by inversion. For a binary substrate ($b_{\max} = 2^w$ for $w$-bit words), the ternary crossover occurs when $b_{\text{eff}} = 3$:
\begin{equation}
    3 = 1 + (2^w - 1)(1 - e^{-\Delta E / \kB T_3})
\end{equation}
\begin{equation}
    e^{-\Delta E / \kB T_3} = 1 - \frac{2}{2^w - 1} \approx 1 - 2^{1-w}
\end{equation}
\begin{equation}
    T_3 \approx \frac{\Delta E}{\kB \cdot 2^{1-w}} = \frac{\Delta E \cdot 2^{w-1}}{\kB}
\end{equation}

For a single bit ($w = 1$): $T_3 = \Delta E / \kB$. This is the temperature at which a single binary state becomes ``ternary''---the two bit states plus the superposition (indistinguishable) state give three effective partition cells.

For $w = 64$ and $\Delta E = 0.1$ eV (typical CMOS threshold):
\begin{equation}
    T_3 \approx \frac{0.1 \times 1.6 \times 10^{-19} \times 2^{63}}{1.38 \times 10^{-23}} \approx 10^{36} \text{ K}
\end{equation}

This absurd temperature confirms that \emph{conventional silicon processors cannot be thermally converted to categorical mode}. The energy spacing $\Delta E$ is too large. Categorical conversion requires a substrate where $\Delta E$ is comparable to $\kB T$---that is, a system operating near its partition extinction threshold.
\end{proof}

\subsection{The Biological Operating Point}

\begin{proposition}[Biological Categorical Conversion]
\label{prop:biological}
Lipid membranes operate naturally at the ternary crossover temperature.
\end{proposition}

\begin{proof}
In a lipid membrane, the relevant states are conformational isomers of the hydrocarbon chains. Each C--C bond has three rotational states (trans, gauche$^+$, gauche$^-$) with energy spacing:
\begin{equation}
    \Delta E_g = E_{\text{gauche}} - E_{\text{trans}} \approx 2.1 \text{ kJ/mol} \approx 3.5 \times 10^{-21} \text{ J}
\end{equation}

The ternary crossover temperature:
\begin{equation}
    T_3 = \frac{\Delta E_g}{\kB \ln(3/2)} = \frac{3.5 \times 10^{-21}}{1.38 \times 10^{-23} \times 0.405} \approx 626 \text{ K}
\end{equation}

But this is for a single bond with $b_{\max} = 3$. For a chain of $n_c = 16$ carbons with 15 rotatable bonds, the effective energy spacing per collective mode is reduced by cooperative effects:
\begin{equation}
    \Delta E_{\text{eff}} \approx \frac{\Delta E_g}{\sqrt{n_c - 1}} \approx \frac{3.5 \times 10^{-21}}{\sqrt{15}} \approx 9.0 \times 10^{-22} \text{ J}
\end{equation}

The collective ternary crossover:
\begin{equation}
    T_3^{\text{collective}} \approx \frac{9.0 \times 10^{-22}}{1.38 \times 10^{-23} \times 0.405} \approx 161 \text{ K}
\end{equation}

Including interchain cooperative corrections (the chain is not isolated but coupled to neighbors through van der Waals interactions):
\begin{equation}
    T_3^{\text{membrane}} \approx T_3^{\text{collective}} \times \frac{E_{\text{vdW}} + \Delta E_g}{\Delta E_g} \approx 161 \times \frac{4.5 + 2.1}{2.1} \approx 161 \times 3.1 \approx 500 \text{ K}
\end{equation}

The range $310$--$500$ K brackets the biological operating temperature of $310$ K. Biological membranes operate at or slightly below $T_3$, in the regime where conformational states are partially distinguishable---neither fully binary nor fully extinct. This is the \emph{categorical computing regime}.
\end{proof}

\begin{remark}[Why $37°$C?]
The biological operating temperature is not arbitrary. It is tuned to place the lipid membrane at the ternary crossover---the regime of maximum categorical computational power. Evolution did not ``choose'' $37°$C for protein stability alone; it chose it because this temperature places the membrane processor at optimal partition resolution.
\end{remark}

\begin{figure*}[!htbp]
    \centering
    \includegraphics[width=0.9\textwidth]{panels/panel_exp10_effective_base.png}
    \caption{\textbf{Effective encoding base $b_{\mathrm{eff}}(T)$ as a continuous function of temperature.}
    (\textbf{A}) $b_{\mathrm{eff}}(T)$ for five physical systems spanning energy gaps from $2\,$kJ/mol (van der Waals) to $106\,$kJ/mol (silicon bandgap). Silicon transistors (blue) remain locked at $b_{\mathrm{eff}} = 3.0$ across the entire temperature range ($\Delta E \gg k_BT$), while lipid chains (red) and weak interactions (gold) show continuous base variation near biological temperatures---the prerequisite for thermal base conversion.
    (\textbf{B}) $b_{\mathrm{eff}}$ for lipid chains at four values of $b_{\max} \in \{2, 3, 5, 10\}$, demonstrating that the functional form $b_{\mathrm{eff}} = 1 + (b_{\max} - 1)(1 - e^{-\Delta E / k_BT})$ smoothly interpolates between unity (partition extinction) and $b_{\max}$ (full resolution) regardless of the maximum base.
    (\textbf{C}) 3D surface of $b_{\mathrm{eff}}$ over the $(T, b_{\max})$ plane for lipid chain energetics. The surface shows a plateau at high $b_{\max}$ and low $T$ (fully resolved partitions) that smoothly descends toward $b_{\mathrm{eff}} = 1$ at high temperatures, with the biological operating point ($T = 310\,$K, $b_{\max} = 3$) situated on the steepest gradient---maximising sensitivity to thermal modulation.
    (\textbf{D}) Partition depth increment $\delta M = 1/\ln b_{\mathrm{eff}}$ for all five systems versus temperature. The divergence of $\delta M$ as $b_{\mathrm{eff}} \to 1$ (partition extinction) for weak interactions demonstrates the fundamental limit where state counting ceases and categorical processing becomes the only viable computational modality.}
    \label{fig:exp10_effective_base}
\end{figure*}

%==============================================================================
\section{The Membrane as Transducer}
\label{sec:membrane}
%==============================================================================

\subsection{Virtual Membranes}

\begin{definition}[Virtual Membrane]
\label{def:virtual_membrane}
A \emph{virtual membrane} is a boundary condition in partition space that separates regions of different effective partition base $b_{\text{eff}}$. It is defined by:
\begin{enumerate}
    \item A partition surface $\Sigma$ in the state space of the system
    \item A partition depth gradient $\nabla\Depth$ across $\Sigma$
    \item A base conversion function $b_{\text{eff}}(\mathbf{x})$ that varies across $\Sigma$
\end{enumerate}
A virtual membrane need not be a physical lipid bilayer---it is any structure that mediates between partition regimes.
\end{definition}

\begin{theorem}[Membrane Transduction Theorem]
\label{thm:transduction}
A virtual membrane separating a binary region ($b = 2$) from a categorical region ($b = b_c$) performs base conversion with information conservation:
\begin{equation}
    I_{\text{binary}} = I_{\text{categorical}}
\end{equation}
\begin{equation}
    n_{\text{bits}} \cdot \log_2 2 = \Depth_{\text{cat}} \cdot \log_2 b_c
\end{equation}
\begin{equation}
    \therefore \Depth_{\text{cat}} = \frac{n_{\text{bits}}}{\log_2 b_c}
\end{equation}
\end{theorem}

\begin{proof}
Information is conserved across the membrane because the membrane is a partition boundary, not an information sink. The total number of distinguishable configurations on each side must match at the boundary.

On the binary side: $2^{n_{\text{bits}}}$ configurations, information $I = n_{\text{bits}}$ bits.

On the categorical side: $b_c^{\Depth_{\text{cat}}}$ configurations, information $I = \Depth_{\text{cat}} \log_2 b_c$ bits.

Equating: $n_{\text{bits}} = \Depth_{\text{cat}} \log_2 b_c$, giving $\Depth_{\text{cat}} = n_{\text{bits}} / \log_2 b_c$.

For binary-to-ternary conversion ($b_c = 3$):
\begin{equation}
    \Depth_{\text{ternary}} = \frac{n_{\text{bits}}}{\log_2 3} = \frac{n_{\text{bits}}}{1.585} \approx 0.631 \cdot n_{\text{bits}}
\end{equation}

A 64-bit binary word converts to a ternary partition depth of $\Depth \approx 40.4$. The information content is identical; the representation is more compact.
\end{proof}

\subsection{The Capacity Formula as Base Converter}

\begin{theorem}[Capacity Formula Conversion]
\label{thm:capacity_conversion}
The capacity formula $C(n) = 2n^2$ provides a natural bijection between binary state counts and ternary partition depth:
\begin{equation}
    n_{\text{binary states}} = C(n) = 2n^2 \quad \Leftrightarrow \quad n = \sqrt{\frac{N_{\text{states}}}{2}}
\end{equation}

The partition depth at level $n$ in ternary encoding is:
\begin{equation}
    \Depth_3(n) = \log_3 C_{\text{tot}}(n) = \log_3\left(\frac{n(n+1)(2n+1)}{3}\right)
\end{equation}

The binary-to-ternary conversion efficiency:
\begin{equation}
    \eta(n) = \frac{\Depth_3(n)}{\log_2 C_{\text{tot}}(n)} = \frac{1}{\log_2 3} = 0.631
\end{equation}
is constant and equals the information-theoretic ratio $\log_3 2$.
\end{theorem}

\begin{proof}
The capacity formula gives the number of states at each shell:
\begin{center}
\begin{tabular}{ccccc}
\toprule
$n$ & $C(n)$ & $C_{\text{tot}}(n)$ & $\log_2 C_{\text{tot}}$ & $\log_3 C_{\text{tot}}$ \\
\midrule
1 & 2 & 2 & 1.00 & 0.63 \\
2 & 8 & 10 & 3.32 & 2.10 \\
3 & 18 & 28 & 4.81 & 3.03 \\
4 & 32 & 60 & 5.91 & 3.73 \\
5 & 50 & 110 & 6.78 & 4.28 \\
10 & 200 & 770 & 9.59 & 6.05 \\
20 & 800 & 5740 & 12.49 & 7.88 \\
\bottomrule
\end{tabular}
\end{center}

The ratio $\log_3 / \log_2 = 1/\log_2 3 = 0.631$ is constant for all $n$. The capacity formula provides a perfect base conversion because it encodes the same partition structure---the angular momentum hierarchy---in both bases.

The factor of 2 in $C(n) = 2n^2$ is the binary element (chirality $s = \pm 1/2$), while $n^2$ is the angular structure $\sum_{\ell=0}^{n-1}(2\ell+1)$, which has natural ternary character (each $\ell$ contributes $2\ell+1$ states, and $2\ell+1 \in \{1, 3, 5, 7, \ldots\}$---odd numbers, with 3 being the first non-trivial case).
\end{proof}

\subsection{The Conversion Protocol}

\begin{definition}[Categorical Conversion Protocol]
\label{def:protocol}
The conversion of a binary computation to categorical form proceeds through:
\begin{enumerate}
    \item \textbf{Encoding}: Map $n$-bit binary input to partition coordinates $(n, \ell, m, s)$ using the state index bijection.
    \item \textbf{Membrane crossing}: The encoded state passes through the virtual membrane, where temperature-mediated partition extinction converts binary distinguishability to categorical distinguishability.
    \item \textbf{Navigation}: The categorical state navigates S-entropy space $(\Sk, \St, \Se)$ toward the solution at the $\eps$-boundary.
    \item \textbf{Decoding}: At the $\eps$-boundary, the categorical state is re-encoded in binary through the inverse bijection.
\end{enumerate}
\end{definition}

\begin{theorem}[Conversion Fidelity]
\label{thm:fidelity}
The categorical conversion protocol preserves information exactly:
\begin{equation}
    H(\text{output}) = H(\text{input})
\end{equation}
where $H$ is Shannon entropy.
\end{theorem}

\begin{proof}
Each step is bijective:
\begin{itemize}
    \item Binary $\to$ partition coordinates: bijection by Theorem~\ref{thm:capacity_conversion}
    \item Membrane crossing: information conservation by Theorem~\ref{thm:transduction}
    \item Navigation: deterministic trajectory in S-entropy space (no information loss)
    \item Categorical $\to$ binary: inverse bijection
\end{itemize}

The composition of bijections is a bijection. No information is created or destroyed.
\end{proof}

\begin{figure*}[!htbp]
    \centering
    \includegraphics[width=0.9\textwidth]{panels/panel_exp11.png}
    \caption{\textbf{Universal capacity conversion ratio and inter-binary state accessibility.}
    (\textbf{A}) Partition depth in binary ($M_2$, blue), ternary ($M_3$, magenta), and effective base at $310\,$K ($M_{\mathrm{eff}}$, green) representations versus partition level $n$. All three grow as $\log C(n)$ with constant multiplicative offsets, confirming base-independent scaling of the capacity formula.
    (\textbf{B}) Number of inter-binary states (states addressable by the effective base but falling between binary grid points) versus $n$, coloured by the effective conversion ratio $R_{\mathrm{eff}}$. The quadratic growth demonstrates that the categorical advantage scales with system size: larger state spaces harbour proportionally more states invisible to binary encoding.
    (\textbf{C}) 3D surface of partition depth $\log_b(2n^2)$ over the $(n, b)$ plane for continuous base $b \in [2, 4]$. The smooth surface confirms that partition depth is an analytic function of both system size and encoding base, with no discontinuities or phase transitions---the mathematical foundation of continuous base conversion.
    (\textbf{D}) Conversion ratios $R_{2\to3} = \ln 2 / \ln 3$ (blue) and $R_{\mathrm{eff}}$ at $310\,$K (red dashed) versus $n$. Both are rigorously constant ($\sigma < 10^{-16}$), establishing the universality theorem: base conversion ratios depend only on the bases, never on the partition level or state count.}
    \label{fig:exp11_conversion}
\end{figure*}

%==============================================================================
\section{The Non-Quantizable Processor}
\label{sec:processor}
%==============================================================================

\subsection{What Makes It Non-Quantizable}

\begin{theorem}[Non-Quantizability]
\label{thm:non_quantizable}
The categorical processor has no fixed computational unit. Its partition depth increment per operation is:
\begin{equation}
    \delta\Depth(T) = \frac{1}{\ln b_{\text{eff}}(T)}
\end{equation}
which is a continuous function of temperature.
\end{theorem}

\begin{proof}
In a binary processor, each operation changes one bit, giving $\delta\Depth = 1 / \ln 2$. This is fixed.

In the categorical processor, the effective base $b_{\text{eff}}(T)$ varies continuously with temperature. The depth increment per operation:
\begin{equation}
    \delta\Depth(T) = \frac{1}{\ln b_{\text{eff}}(T)} = \frac{1}{\ln[1 + (b_{\max}-1)(1-e^{-\Delta E/\kB T})]}
\end{equation}

This is a smooth function of $T$ with no discrete steps. At no temperature is there a ``natural unit'' of computation. The processor's granularity is continuously tunable:

\begin{center}
\begin{tabular}{lccc}
\toprule
Regime & $T/T_c$ & $b_{\text{eff}}$ & $\delta\Depth$ \\
\midrule
Deep binary & $\gg 1$ & $b_{\max}$ & $1/\ln b_{\max}$ \\
Binary & $\sim 10$ & 2 & 1.443 \\
Ternary & $\sim 3$ & 3 & 0.910 \\
Pentary & $\sim 2$ & 5 & 0.621 \\
High base & $\sim 1.5$ & 10 & 0.434 \\
Near extinction & $\sim 1.01$ & 100 & 0.217 \\
Continuous & $\to 1$ & $\to \infty$ & $\to 0$ \\
\bottomrule
\end{tabular}
\end{center}

As $b_{\text{eff}} \to \infty$, $\delta\Depth \to 0$---the processor achieves \emph{infinitesimal} partition depth resolution. This is the continuous limit: computation without quantization.
\end{proof}

\begin{corollary}[No Computational Atom]
\label{cor:no_atom}
Unlike binary processors (where the bit is the indivisible unit) or quantum processors (where the qubit is the unit), the categorical processor has no irreducible computational element. Any apparent ``unit'' is an artifact of the current temperature, not a property of the architecture.
\end{corollary}

\subsection{Processor Architecture}

\begin{definition}[Categorical Processor]
\label{def:cat_processor}
A categorical processor $\mathcal{P}$ is defined by the tuple:
\begin{equation}
    \mathcal{P} = (\Sigma, T(\mathbf{x}), b_{\text{eff}}(\mathbf{x}), \omega(\mathbf{x}))
\end{equation}
where:
\begin{itemize}
    \item $\Sigma$: The membrane surface (virtual or physical)
    \item $T(\mathbf{x})$: Temperature field across $\Sigma$
    \item $b_{\text{eff}}(\mathbf{x})$: Effective base field (derived from $T$)
    \item $\omega(\mathbf{x})$: Oscillation frequency field (processing rate)
\end{itemize}
\end{definition}

\begin{theorem}[Processing Rate]
\label{thm:processing_rate}
The categorical processing rate of a membrane of area $A$ is:
\begin{equation}
    R_{\text{cat}} = \frac{2}{A_L} \int_A \frac{\omega(\mathbf{x})}{2\pi} \cdot \Depth_{\text{mode}}(\mathbf{x}) \, dA
\end{equation}
where $A_L$ is the area per processing element and $\Depth_{\text{mode}}$ is the partition depth per oscillation cycle.
\end{theorem}

\begin{proof}
Each processing element (lipid, virtual oscillator, or other partition unit) at position $\mathbf{x}$ on the membrane contributes:
\begin{equation}
    r(\mathbf{x}) = \frac{\omega(\mathbf{x})}{2\pi} \cdot \Depth_{\text{mode}}(\mathbf{x})
\end{equation}
partition depth per second. The factor of 2 accounts for both sides of the membrane (both leaflets in a physical bilayer).

For a uniform membrane with $\omega = 10^{11}$ rad/s (chain isomerization), $\Depth_{\text{mode}} = \log_2 3 \approx 1.585$ (ternary conformational partition), and $A_L = 0.64$ nm$^2$:
\begin{equation}
    R_{\text{cat}} = \frac{2 \times A}{0.64 \times 10^{-18}} \times \frac{10^{11}}{2\pi} \times 1.585
\end{equation}

For $A = 1000$ $\mu$m$^2 = 10^{-9}$ m$^2$:
\begin{equation}
    R_{\text{cat}} = \frac{2 \times 10^{-9}}{6.4 \times 10^{-19}} \times 1.59 \times 10^{10} \times 1.585 \approx 7.9 \times 10^{19} \approx 10^{20} \text{ ops/s}
\end{equation}
\end{proof}

\subsection{The S-Entropy State Space}

\begin{definition}[S-Entropy Coordinates]
\label{def:s_entropy}
The categorical processor operates in S-entropy coordinate space $\Scoord = (\Sk, \St, \Se) \in [0,1]^3$:
\begin{align}
    \Sk &= -\sum_i p_i^{(k)} \ln p_i^{(k)} / \ln N_k & \text{(knowledge entropy)} \\
    \St &= -\sum_i p_i^{(t)} \ln p_i^{(t)} / \ln N_t & \text{(temporal entropy)} \\
    \Se &= -\sum_i p_i^{(e)} \ln p_i^{(e)} / \ln N_e & \text{(evolution entropy)}
\end{align}
where $p_i^{(\alpha)}$ are the probability distributions over the respective partition coordinates and $N_\alpha$ are normalizing factors ensuring $S_\alpha \in [0,1]$.
\end{definition}

\begin{theorem}[S-Entropy Navigation]
\label{thm:navigation}
Computation in the categorical processor is navigation through S-entropy space. The computational state $\Scoord(t)$ evolves according to:
\begin{equation}
    \frac{d\Scoord}{dt} = -\nabla_{\Scoord} \mathcal{F}(\Scoord) + \boldsymbol{\xi}(t)
\end{equation}
where $\mathcal{F}(\Scoord)$ is the categorical free energy landscape encoding the problem, and $\boldsymbol{\xi}(t)$ is thermal noise with $\langle \xi_\alpha(t) \xi_\beta(t') \rangle = 2D_\alpha \delta_{\alpha\beta}\delta(t-t')$.
\end{theorem}

\begin{proof}
The categorical state $\Scoord$ is a point in the bounded cube $[0,1]^3$. The problem to be solved defines a free energy landscape $\mathcal{F}(\Scoord)$ with minima at solution states and barriers between them.

The membrane's oscillatory dynamics generate Brownian motion in S-space with diffusion coefficient:
\begin{equation}
    D_\alpha = \frac{\kB T}{6\pi\eta R_{\text{eff}}}
\end{equation}
where $\eta$ is an effective viscosity in partition space and $R_{\text{eff}}$ is the effective size of the computational state.

The drift term $-\nabla_{\Scoord}\mathcal{F}$ drives the system toward low free energy (solution states). The noise term $\boldsymbol{\xi}$ enables escape from local minima---this is the categorical analog of simulated annealing, but performed physically by the membrane's thermal fluctuations.

The solution is found when $\Scoord$ reaches the $\eps$-boundary of the target region:
\begin{equation}
    |\Scoord - \Scoord^*| < \eps
\end{equation}
where $\Scoord^*$ is the solution and $\eps$ is the precision limit set by the G\"odelian residue $\Gres$.
\end{proof}

\begin{figure*}[!htbp]
    \centering
    \includegraphics[width=0.9\textwidth]{panels/panel_exp14.png}
    \caption{\textbf{S-entropy navigation in $(S_k, S_t, S_e)$ coordinate space.}
    (\textbf{A}) Grouped bar chart comparing search steps required by linear scan (red), binary search (blue), and categorical search (green) across problem sizes $N \in \{100, 500, 1000, 5000, 10000\}$ on a logarithmic scale. Categorical search consistently requires $2$--$3$ steps regardless of $N$, compared to $7$--$14$ for binary search, demonstrating the power of simultaneous 3D coordinate refinement.
    (\textbf{B}) Euclidean distance to target versus navigation step along the categorical search trajectory. Convergence is achieved within $2$ steps, with distance dropping from $\sim 0.5$ to below $10^{-6}$, confirming the exponential convergence rate $d(k) \sim b_{\mathrm{eff}}^{-k}$ predicted by the partition refinement theorem.
    (\textbf{C}) 3D trajectory through $(S_k, S_t, S_e)$ entropy space, showing the categorical search path from the initial estimate (centre of the unit cube) to the target state. Each waypoint (coloured by step number) represents one partition refinement, with the resolution shrinking by a factor of $b_{\mathrm{eff}}^{-1} \approx 0.37$ per step in each dimension simultaneously.
    (\textbf{D}) Theoretical categorical step count $\log_{b_{\mathrm{eff}}}(N)$ (green line) versus actual measured steps (red markers) for each problem size. The close agreement validates the $O(\log_{b_{\mathrm{eff}}} N)$ complexity bound and confirms that the hierarchical partition structure enables near-optimal navigation without index structures.}
    \label{fig:exp14_navigation}
\end{figure*}

%==============================================================================
\section{Categorical Thermodynamics}
\label{sec:thermodynamics}
%==============================================================================

\subsection{Entropy of Conversion}

\begin{theorem}[Conversion Entropy]
\label{thm:conversion_entropy}
The entropy change during conversion from binary (base 2) to categorical (base $b_c$) encoding of $N$ states is:
\begin{equation}
    \Delta S_{\text{conversion}} = N \kB (\ln b_c - \ln 2)
\end{equation}
For ternary ($b_c = 3$): $\Delta S = N \kB \ln(3/2) > 0$.
\end{theorem}

\begin{proof}
In binary encoding, $N$ states occupy $\lceil \log_2 N \rceil$ bits, each with depth $\Depth = 1$. The entropy per state:
\begin{equation}
    s_{\text{binary}} = \kB \ln 2
\end{equation}

In base-$b_c$ encoding, the same $N$ states occupy $\lceil \log_{b_c} N \rceil$ digits, each with depth $\Depth = 1$. The entropy per state:
\begin{equation}
    s_{\text{cat}} = \kB \ln b_c
\end{equation}

The entropy change:
\begin{equation}
    \Delta S = N(s_{\text{cat}} - s_{\text{binary}}) = N\kB(\ln b_c - \ln 2)
\end{equation}

For $b_c > 2$: $\Delta S > 0$. The conversion to higher base \emph{increases} entropy. This is thermodynamically spontaneous at all temperatures---the system naturally evolves toward higher-base encoding.
\end{proof}

\begin{corollary}[Spontaneous Categorification]
\label{cor:spontaneous}
The conversion from binary to categorical encoding is thermodynamically spontaneous ($\Delta G < 0$ at constant $T$, $P$). Binary encoding is a metastable state maintained by energy barriers; categorical encoding is the equilibrium.
\end{corollary}

\subsection{Carnot Efficiency of Conversion}

\begin{theorem}[Categorical Carnot Bound]
\label{thm:carnot}
The efficiency of categorical conversion---the fraction of input computational work recovered as categorical output---is bounded by:
\begin{equation}
    \eta_{\text{cat}} \leq 1 - \frac{T_{\text{cold}}}{T_{\text{hot}}} = 1 - \frac{T_{\text{extinction}}}{T_{\text{binary}}}
\end{equation}
where $T_{\text{binary}}$ is the temperature of the binary substrate and $T_{\text{extinction}}$ is the partition extinction temperature.
\end{theorem}

\begin{proof}
The categorical converter is a heat engine operating between two temperatures:
\begin{itemize}
    \item $T_{\text{hot}} = T_{\text{binary}}$: The temperature at which binary states are fully distinguishable
    \item $T_{\text{cold}} = T_{\text{extinction}}$: The temperature at which partition extinction occurs
\end{itemize}

The converter extracts work from the temperature difference by converting high-entropy binary encoding to low-entropy categorical encoding. By the second law, the maximum efficiency is the Carnot efficiency:
\begin{equation}
    \eta_{\max} = 1 - \frac{T_{\text{extinction}}}{T_{\text{binary}}}
\end{equation}

For a biological membrane operating between $T_{\text{binary}} \approx 500$ K (fully fluid) and $T_{\text{extinction}} \approx 270$ K (gel phase):
\begin{equation}
    \eta_{\max} = 1 - \frac{270}{500} = 0.46
\end{equation}

At the biological operating point $T = 310$ K:
\begin{equation}
    \eta_{\text{bio}} \approx 0.38
\end{equation}

This is remarkably close to the actual efficiency of ATP synthesis ($\sim 40\%$), suggesting that biological energy conversion is operating near the categorical Carnot limit.
\end{proof}

\subsection{Irreversibility and the Arrow of Computation}

\begin{theorem}[Computational Arrow]
\label{thm:arrow}
Categorical computation has an intrinsic direction: from higher to lower partition depth. This direction is the \emph{arrow of computation}, which is equivalent to the thermodynamic arrow of time.
\end{theorem}

\begin{proof}
Each partition operation generates entropy $\Delta S > 0$ (by the categorical second law). The cumulative entropy:
\begin{equation}
    S(t) = S(0) + \int_0^t \dot{S}(t') \, dt' > S(0)
\end{equation}

The partition depth of the system decreases as states merge:
\begin{equation}
    \Depth(t) \leq \Depth(0)
\end{equation}

Computation proceeds from high partition depth (many distinguishable states, high information) to low partition depth (fewer states, solution found). This is irreversible: you cannot ``un-compute'' without paying at least $T\Delta S$ in work.

The arrow of computation IS the arrow of time, viewed through the partition lens. The passage of time is the accumulation of partition transitions; computation is the directed accumulation toward a solution.
\end{proof}

\begin{figure*}[!htbp]
    \centering
    \includegraphics[width=0.9\textwidth]{panels/panel_exp13.png}
    \caption{\textbf{Thermodynamics of binary-to-categorical conversion.}
    (\textbf{A}) Stacked area plot of Carnot efficiency $\eta_C = 1 - T_{\mathrm{env}}/T$ (red) and partition efficiency $\eta_P = 1 - \ln 2 / \ln b_{\mathrm{eff}}$ (blue) versus temperature. At $T = 310\,$K, the partition contribution ($30.5\%$) dominates over the thermal contribution ($3.8\%$), yielding a combined efficiency of $33.2\%$---the thermodynamic cost of converting binary to categorical encoding.
    (\textbf{B}) Number of base-conversion events sustainable per ATP hydrolysis ($\Delta G_{\mathrm{ATP}} = 30.5\,$kJ/mol) versus temperature. At body temperature, approximately $39$ conversion events can be powered by a single ATP, suggesting that biological categorical processing is energetically economical.
    (\textbf{C}) 3D surface of combined conversion efficiency over the $(T, \Delta E)$ plane, where $\Delta E$ spans the range of biologically relevant energy gaps ($0.5$--$5.0\,$kJ/mol). The surface reveals an efficiency plateau above $\Delta E > 3\,$kJ/mol, indicating that lipid chain barriers ($\sim 5\,$kJ/mol) operate in the saturated regime where further barrier increases provide diminishing returns.
    (\textbf{D}) Work extracted per conversion event $W = k_BT\ln(b_{\mathrm{eff}}/2)$ (blue) and mixing free energy cost $\Delta G_{\mathrm{mix}}$ (red dashed) versus temperature. The positive $W$ above $\sim 290\,$K confirms that binary-to-categorical conversion is thermodynamically spontaneous at biological temperatures, with the driving force increasing linearly with $T$.}
    \label{fig:exp13_thermo}
\end{figure*}

%==============================================================================
\section{The Categorical Speedup}
\label{sec:speedup}
%==============================================================================

\subsection{Binary vs. Categorical Complexity}

\begin{theorem}[Categorical Speedup Theorem]
\label{thm:speedup}
For a search problem in a space of $N$ states:
\begin{enumerate}
    \item Binary sequential search: $O(N)$ partition transitions
    \item Binary parallel search: $O(N/p)$ transitions with $p$ processors
    \item Categorical navigation: $O(\log N)$ partition transitions
\end{enumerate}
The categorical speedup is exponential: $N / \log N$.
\end{theorem}

\begin{proof}
\textbf{Binary sequential}: Each state must be individually examined. With $N$ states and 1 transition per examination: $O(N)$.

\textbf{Binary parallel}: With $p$ processors, each examines $N/p$ states: $O(N/p)$. Even with $p = N$ processors (maximum parallelism), the complexity is $O(1)$ per processor but $O(N)$ total transitions.

\textbf{Categorical navigation}: The solution exists at a point $\Scoord^*$ in S-entropy space $[0,1]^3$. The problem is encoded as a free energy landscape $\mathcal{F}(\Scoord)$ with a minimum at $\Scoord^*$.

Navigation proceeds by halving the search volume at each step. From any point $\Scoord$, the gradient $\nabla\mathcal{F}$ indicates which half of the remaining volume contains the solution. Each gradient evaluation is one categorical operation (one oscillation cycle).

Starting from the full volume $V_0 = 1$ (unit cube), after $k$ steps the remaining volume is:
\begin{equation}
    V_k = V_0 / 2^k
\end{equation}

The solution is found when $V_k < \eps^3$ (the $\eps$-boundary volume):
\begin{equation}
    k = \log_2(1/\eps^3) = 3\log_2(1/\eps)
\end{equation}

For $N$ states uniformly distributed in $[0,1]^3$, the inter-state spacing is $\eps \sim N^{-1/3}$:
\begin{equation}
    k = 3\log_2(N^{1/3}) = \log_2 N
\end{equation}

Therefore: categorical search complexity is $O(\log N)$.

The speedup over binary sequential: $N / \log N$.
For $N = 10^6$: speedup $= 10^6 / 20 = 50{,}000\times$.
For $N = 10^{12}$: speedup $= 10^{12} / 40 = 2.5 \times 10^{10}\times$.
\end{proof}

\subsection{Why Categorical Navigation Achieves Logarithmic Complexity}

\begin{proposition}[Gradient Availability]
\label{prop:gradient}
In S-entropy space, the free energy gradient $\nabla\mathcal{F}$ is always available as a physical observable---it is encoded in the membrane's oscillatory dynamics.
\end{proposition}

\begin{proof}
The membrane's oscillatory state encodes $\Scoord$ (knowledge, temporal, and evolution entropy are derived from oscillation amplitudes, phases, and frequencies respectively). The gradient $\nabla\mathcal{F}$ corresponds to the direction in which oscillatory modes are most strongly damped.

This gradient is not computed---it is \emph{measured} by the membrane's physical dynamics. The membrane naturally evolves in the direction of decreasing $\mathcal{F}$ because:
\begin{enumerate}
    \item Oscillatory modes with lower free energy have longer lifetimes
    \item Higher-energy modes decay through partition transitions
    \item The surviving mode distribution encodes the gradient
\end{enumerate}

Each oscillation cycle samples the local free energy landscape and takes one step in the gradient direction. This is $O(1)$ per step, and $O(\log N)$ steps suffice to navigate to the solution.
\end{proof}

\subsection{The $\eps$-Boundary and Computational Completeness}

\begin{theorem}[$\eps$-Boundary Precision]
\label{thm:eps_boundary}
The categorical computation reaches the $\eps$-boundary but never achieves exact closure. The precision is:
\begin{equation}
    |\Scoord - \Scoord^*| \geq \eps = \frac{\Gres}{N^{1/3}}
\end{equation}
where $\Gres$ is the G\"odelian residue---the irreducible imprecision arising from the impossibility of complete self-reference in bounded systems.
\end{theorem}

\begin{proof}
The G\"odelian residue arises from the partition framework's equivalent of G\"odel's incompleteness theorem: a bounded system cannot fully specify its own partition structure because the specification itself requires partition operations that alter the structure being specified.

The residue is:
\begin{equation}
    \Gres = \frac{\hbar}{2\kB T \tau_p}
\end{equation}
where $\tau_p$ is the partition operation time and $\hbar/(2\kB T)$ is the quantum thermal time.

At biological temperature ($T = 310$ K) with $\tau_p \sim 10^{-11}$ s:
\begin{equation}
    \Gres \approx \frac{1.05 \times 10^{-34}}{2 \times 1.38 \times 10^{-23} \times 310 \times 10^{-11}} \approx 1.2 \times 10^{-3}
\end{equation}

The $\eps$-boundary precision for $N = 10^6$:
\begin{equation}
    \eps = \frac{1.2 \times 10^{-3}}{(10^6)^{1/3}} = \frac{1.2 \times 10^{-3}}{100} = 1.2 \times 10^{-5}
\end{equation}

This corresponds to $\sim 17$ bits of precision---sufficient for most computational tasks. Higher precision requires either lower temperature (increasing $\tau_p$) or larger systems (increasing $N$).
\end{proof}

\begin{figure*}[!htbp]
    \centering
    \includegraphics[width=0.9\textwidth]{panels/panel_exp12.png}
    \caption{\textbf{Categorical speedup: $O(\log N)$ versus $O(N)$ state-space search.}
    (\textbf{A}) Log-log plot of operations required versus problem size $N$ for linear scan (blue), binary search (magenta), and categorical search at $b_{\mathrm{eff}} = 2.71$ (green). The categorical curve lies strictly below binary search, with the gap widening as $N$ increases, reflecting the constant factor $\log_2(b_{\mathrm{eff}}) / 1 = 1.44$ speedup per step.
    (\textbf{B}) Speedup factor (linear operations / categorical operations) versus $\log_{10} N$. The speedup grows super-linearly with problem size, reaching $\sim 10^{10}$ for $N = 10^{12}$, demonstrating that categorical processing provides its greatest advantage for large-scale biological state spaces.
    (\textbf{C}) 3D surface of search operations $\log_b(N)$ over the $(\log_{10} N, b)$ plane for continuous base $b \in [2, 4]$. The surface monotonically decreases with increasing base, confirming that higher effective bases yield fewer search steps for any problem size, with the steepest gains in the biologically relevant regime $b \in [2.5, 3.0]$.
    (\textbf{D}) Time speedup (binary wall-clock / categorical wall-clock) versus problem size, accounting for the $10\times$ clock advantage of lipid isomerisation ($\sim 10^{10}\,$Hz) over silicon switching ($\sim 10^9\,$Hz). The combined algorithmic and clock-rate advantage yields a $\sim 14\times$ speedup factor that is essentially constant across problem sizes.}
    \label{fig:exp12_speedup}
\end{figure*}

%==============================================================================
\section{Biological Implementation}
\label{sec:biological}
%==============================================================================

\subsection{The Cell Membrane as Categorical Processor}

\begin{theorem}[Biological Categorical Processing]
\label{thm:bio_processing}
A eukaryotic cell membrane of area $A \approx 1000$ $\mu$m$^2$ implements a categorical processor with:
\begin{enumerate}
    \item Processing rate: $R \approx 10^{20}$ ops/s
    \item Effective base: $b_{\text{eff}} \approx 2.5$--$3$ at $T = 310$ K
    \item Partition depth per cycle: $\Depth \approx 1.3$--$1.6$
    \item S-entropy resolution: $\delta S \approx 10^{-3}$
    \item Computational precision: $\sim 17$ bits equivalent
\end{enumerate}
\end{theorem}

\begin{proof}
From the membrane biophysics established in the companion paper:
\begin{itemize}
    \item Lipid count: $N_L = 2A/A_L \approx 3.1 \times 10^9$
    \item Conformational frequency: $\omega \sim 10^{11}$ rad/s per lipid
    \item Modes per lipid: $n_c - 1 \approx 15$ rotatable bonds
    \item States per bond: 3 (trans, $g^+$, $g^-$)
\end{itemize}

The effective base at 310 K with $\Delta E_g = 2.1$ kJ/mol:
\begin{equation}
    b_{\text{eff}} = 1 + 2(1 - e^{-\Delta E_g / R T}) = 1 + 2(1 - e^{-2100/2577}) = 1 + 2(1 - 0.444) = 2.11
\end{equation}

Including collective effects that enhance the effective temperature: $b_{\text{eff}} \approx 2.5$--$3.0$.

Processing rate (from Theorem~\ref{thm:processing_rate}):
\begin{equation}
    R = N_L \times 15 \times \frac{10^{11}}{2\pi} \times \log_{b_{\text{eff}}} 3 \approx 3.1 \times 10^9 \times 15 \times 1.6 \times 10^{10} \times 1.2 \approx 8.9 \times 10^{20}
\end{equation}

This exceeds the world's fastest supercomputer ($\sim 10^{18}$ FLOPS) by three orders of magnitude.
\end{proof}

\subsection{Lipid Rafts as Computational Domains}

\begin{theorem}[Raft Computing]
\label{thm:raft_computing}
Lipid rafts---cholesterol-enriched domains in the liquid-ordered phase---function as \emph{categorical coprocessors} with enhanced precision at the cost of reduced processing rate.
\end{theorem}

\begin{proof}
In the raft domain ($L_o$ phase):
\begin{itemize}
    \item Conformational partition depth is reduced (chains ordered by cholesterol)
    \item $b_{\text{eff}}$ is lower than in the surrounding $L_d$ phase
    \item Processing rate decreases: fewer accessible states per cycle
    \item Precision increases: lower noise, more coherent dynamics
\end{itemize}

The raft is a region of \emph{partial partition extinction}: conformational partition operations are suppressed while translational operations persist. This creates a computational domain with:
\begin{equation}
    R_{\text{raft}} < R_{\text{bulk}} \quad \text{but} \quad \eps_{\text{raft}} < \eps_{\text{bulk}}
\end{equation}

The raft trades speed for accuracy. This is the physical basis of signal transduction: signaling proteins concentrate in rafts because the higher precision enables reliable detection of weak signals \cite{simons1997,pike2006}.
\end{proof}

\subsection{Ion Channels as Categorical I/O}

\begin{proposition}[Channel as Categorical Interface]
\label{prop:channel_io}
Ion channels function as input/output ports for the categorical processor, converting ionic concentration gradients (analog signals) into partition state changes (categorical operations).
\end{proposition}

An open ion channel allows specific ions to cross the membrane, changing the partition state of the interior relative to the exterior. The selectivity filter implements a partition coordinate matcher (matching ion partition coordinates to channel coordinates). The gating mechanism implements a categorical switch between ``conduit defined'' and ``conduit undefined'' states.

The channel's conductance $g$ is a partition transport coefficient:
\begin{equation}
    g = \frac{N_{\text{channels}} \cdot e^2}{A \cdot \kB T \cdot \tau_{\text{transit}}}
\end{equation}
where $\tau_{\text{transit}}$ is the ion transit time through the channel.

\subsection{Action Potentials as Categorical Computation Cycles}

\begin{theorem}[Action Potential as Categorical Cycle]
\label{thm:action_potential}
The neuronal action potential is a complete categorical computation cycle:
\begin{enumerate}
    \item \textbf{Input}: Synaptic signals change local membrane composition
    \item \textbf{Threshold}: The categorical free energy landscape $\mathcal{F}(\Scoord)$ crosses a bifurcation
    \item \textbf{Navigation}: The membrane potential traverses S-entropy space from resting to peak
    \item \textbf{Output}: The completed trajectory propagates down the axon
    \item \textbf{Reset}: The refractory period restores the initial S-entropy coordinates
\end{enumerate}
\end{theorem}

\begin{proof}
The action potential proceeds through distinct partition phases:

\textbf{Depolarization}: Na$^+$ channels open, increasing $\Sk$ (the knowledge entropy increases as internal Na$^+$ concentration changes). This is a partition creation event---new ionic distinctions are created.

\textbf{Peak}: The system reaches the $\eps$-boundary of the depolarized state. Maximum $\Sk$.

\textbf{Repolarization}: K$^+$ channels open, Na$^+$ channels close. $\Sk$ decreases. Partition operations reverse direction.

\textbf{Hyperpolarization}: K$^+$ overshoot creates a region of temporary partition extinction (refractory period). $b_{\text{eff}}$ drops below the critical threshold, preventing re-triggering.

\textbf{Recovery}: Ion pumps restore resting concentrations. The system returns to initial $\Scoord$.

The total cycle time ($\sim 2$ ms) and the number of partition transitions ($\sim 10^8$) give a categorical depth per action potential:
\begin{equation}
    \Depth_{\text{AP}} \approx 10^8 \times 1.5 \approx 1.5 \times 10^8
\end{equation}

This is the computational content of a single neuronal firing event---$\sim 10^8$ partition transitions, or approximately $5 \times 10^7$ bits equivalent.
\end{proof}

%==============================================================================
\section{Experimental Predictions}
\label{sec:predictions}
%==============================================================================

\subsection{Temperature-Dependent Computational Resolution}

\begin{enumerate}
    \item \textbf{Prediction}: The effective computational base of a lipid membrane varies continuously with temperature according to:
    \begin{equation}
        b_{\text{eff}}(T) = 1 + 2(1 - e^{-\Delta E_g / \kB T})
    \end{equation}

    \textbf{Test}: Measure conformational dynamics of lipid chains (e.g., by $^2$H NMR order parameters \cite{seelig1974}) as a function of temperature. The order parameter $S_{\text{CD}}$ maps to $b_{\text{eff}}$ through:
    \begin{equation}
        b_{\text{eff}} = \frac{3}{1 + 2|S_{\text{CD}}|}
    \end{equation}

    At $T = 310$ K: $|S_{\text{CD}}| \approx 0.2$, giving $b_{\text{eff}} \approx 2.14$.
    At $T = 350$ K: $|S_{\text{CD}}| \approx 0.15$, giving $b_{\text{eff}} \approx 2.31$.
    At $T = 273$ K: $|S_{\text{CD}}| \approx 0.35$, giving $b_{\text{eff}} \approx 1.76$.
\end{enumerate}

\subsection{Membrane-Mediated Error Correction}

\begin{enumerate}
    \setcounter{enumi}{1}
    \item \textbf{Prediction}: Lipid rafts provide categorical error correction by maintaining higher precision ($\eps_{\text{raft}} < \eps_{\text{bulk}}$) through partial partition extinction.

    \textbf{Test}: Signal transduction fidelity should be measurably higher for receptors in raft domains compared to bulk membrane. The predicted fidelity ratio:
    \begin{equation}
        \frac{F_{\text{raft}}}{F_{\text{bulk}}} = \frac{\eps_{\text{bulk}}}{\eps_{\text{raft}}} = \frac{b_{\text{eff,bulk}}}{b_{\text{eff,raft}}} \approx \frac{2.5}{1.5} \approx 1.7
    \end{equation}

    Raft-associated signaling should show $\sim 70\%$ higher fidelity.
\end{enumerate}

\subsection{Categorical Speedup in Biological Systems}

\begin{enumerate}
    \setcounter{enumi}{2}
    \item \textbf{Prediction}: Biological computations (protein folding, immune recognition, neural pattern completion) achieve $O(\log N)$ scaling where silicon implementations require $O(N)$.

    \textbf{Test}: Compare the time for immune T-cell receptor recognition of an antigen among $N$ possible antigens with the time for a binary search algorithm to identify the antigen from the same library.

    For $N = 10^7$ possible antigens:
    \begin{itemize}
        \item Binary search: $\log_2(10^7) \approx 23$ comparison steps $\times$ $\sim 1$ $\mu$s per step $= 23$ $\mu$s
        \item T-cell recognition: $\sim 1$--$10$ s (experimental \cite{huppa2010})
    \end{itemize}

    The T-cell is slower per step but explores a far richer state space. The categorical prediction is that T-cell recognition time scales as $\log N$:
    \begin{equation}
        t_{\text{recognition}} \propto \log N \cdot \tau_{\text{membrane}}
    \end{equation}
    where $\tau_{\text{membrane}} \sim 1$ s is the membrane equilibration time. This can be tested by measuring recognition time as a function of antigen library size.
\end{enumerate}

\subsection{Conversion Efficiency}

\begin{enumerate}
    \setcounter{enumi}{3}
    \item \textbf{Prediction}: The efficiency of biological energy conversion (ATP synthesis) approaches the categorical Carnot limit:
    \begin{equation}
        \eta_{\text{ATP}} \leq 1 - \frac{T_{\text{gel}}}{T_{\text{fluid}}} \approx 1 - \frac{270}{500} = 0.46
    \end{equation}

    Experimental ATP synthesis efficiency: $\sim 40\%$ \cite{nicholls2013}, consistent with operating near the categorical Carnot bound.
\end{enumerate}

\begin{figure*}[!htbp]
    \centering
    \includegraphics[width=0.9\textwidth]{panels/panel_exp15.png}
    \caption{\textbf{Generalised Landauer limit for arbitrary encoding base.}
    (\textbf{A}) Minimum erasure energy $E_{\min} = k_BT\ln b$ versus temperature for binary ($b = 2$, blue), ternary ($b = 3$, magenta), and categorical ($b = b_{\mathrm{eff}}(T)$, green) encoding. The categorical curve lies between binary and ternary, tracking $b_{\mathrm{eff}}(T)$ continuously. All three scale linearly with $T$, but with distinct slopes set by $\ln b$.
    (\textbf{B}) Maximum operations achievable at a fixed power budget of $1\,$pW for binary (blue) and categorical (red) processors versus temperature on a semi-logarithmic scale. Binary processors sustain more operations per watt due to their lower per-operation energy cost, but each categorical operation processes $\ln b_{\mathrm{eff}} / \ln 2 \approx 1.44$ times more information.
    (\textbf{C}) 3D surface of the generalised Landauer energy $E = k_BT\ln b$ over the $(T, b)$ plane for $b \in [1.5, 10]$. The surface establishes the fundamental thermodynamic floor for computation at any base, revealing that the energy cost of erasure grows logarithmically with base while information capacity grows linearly---the physical basis for the optimal base $b = e \approx 2.718$.
    (\textbf{D}) Entropy production per operation $\Delta S = k_B\ln b$ (in J$\,$mol$^{-1}\,$K$^{-1}$) for binary (blue) and categorical (red dashed) systems versus temperature. The categorical entropy production exceeds binary by $\ln b_{\mathrm{eff}} / \ln 2$, but since each categorical operation resolves more uncertainty, the entropy produced per \emph{bit of resolved information} is identical---confirming that the generalised Landauer principle is base-invariant when measured in natural units.}
    \label{fig:exp15_landauer}
\end{figure*}

%==============================================================================
\section{Discussion}
\label{sec:discussion}
%==============================================================================

\subsection{Unification of Digital and Analog}

The categorical converter resolves a long-standing dichotomy in computation. Digital and analog computing are not distinct paradigms---they are limiting cases of partition-depth computation:

\begin{center}
\begin{tabular}{lccc}
\toprule
Property & Digital & Categorical & Analog \\
\midrule
Base $b$ & 2 & $b_{\text{eff}}(T)$ & $\infty$ \\
Resolution & $\Delta\Depth = 1$ & $\delta\Depth(T)$ & $\delta\Depth \to 0$ \\
Noise sensitivity & Low & Moderate & High \\
Precision per step & 1 bit & $\log_2 b_{\text{eff}}$ bits & $\infty$ (ideal) \\
Scaling & $O(N)$ & $O(\log N)$ & Undefined \\
Physical substrate & Silicon & Membranes & Continuous fields \\
\bottomrule
\end{tabular}
\end{center}

The categorical regime occupies the sweet spot: sufficient resolution for precise computation ($\sim 1.5$ bits/step), low enough noise for reliable operation, and logarithmic scaling through S-entropy navigation.

\subsection{Why Biology Uses Membranes for Computation}

The framework explains why biological computation is membrane-based:

\begin{enumerate}
    \item \textbf{Temperature tuning}: Membranes operate at the ternary crossover, where categorical computation is most powerful.

    \item \textbf{Massive parallelism}: $\sim 10^9$ lipid oscillators per cell provide $\sim 10^{20}$ ops/s.

    \item \textbf{Adaptive resolution}: Temperature fluctuations dynamically adjust $b_{\text{eff}}$, enabling the cell to trade precision for speed.

    \item \textbf{Domain specialization}: Lipid rafts provide high-precision coprocessors for critical computations (signal transduction).

    \item \textbf{I/O integration}: Ion channels provide natural interfaces between categorical computation and ionic signaling.
\end{enumerate}

\subsection{Implications for Artificial Intelligence}

The categorical converter suggests a path beyond both digital and analog AI:

\textbf{Neuromorphic computing}: Current neuromorphic chips use digital or analog circuits to simulate neural dynamics. The categorical framework suggests that the key feature to replicate is not the neuron but the \emph{membrane}---specifically, its role as a categorical converter operating at the ternary crossover.

\textbf{Membrane-inspired architectures}: A processor built on physical oscillators with tunable partition spacing could achieve categorical speedup. Candidate substrates include liquid crystals, colloidal membranes, and soft-matter systems operating near phase transitions \cite{whitesides2002}.

\textbf{The categorical advantage}: For problems that map naturally to S-entropy navigation (pattern recognition, constraint satisfaction, optimization), categorical processors could achieve exponential speedup over binary processors---not through quantum mechanics, but through the thermodynamics of partition extinction.

\subsection{Relationship to Quantum Computing}

The categorical processor is not a quantum computer. The speedup does not arise from superposition or entanglement but from:
\begin{enumerate}
    \item Continuous rather than discrete state representation
    \item S-entropy navigation rather than sequential search
    \item Physical gradient computation rather than algorithmic evaluation
\end{enumerate}

However, the categorical framework is \emph{compatible} with quantum computing. A quantum categorical processor would combine:
\begin{itemize}
    \item Quantum superposition (exploring multiple S-entropy paths simultaneously)
    \item Categorical navigation (logarithmic scaling in S-entropy space)
    \item Membrane transduction (converting between quantum and classical categorical states)
\end{itemize}

The potential speedup: $O(\sqrt{\log N})$---sub-logarithmic complexity.

\subsection{The Discretization Paradox Resolved}

Computational physics faces a persistent paradox: nature is continuous, but we simulate it with discrete bits. The discretization error is unavoidable in binary computation.

The categorical framework resolves this: nature computes categorically (with $b \to \infty$), and binary computation is a lossy compression of the categorical dynamics. The ``discretization error'' is the quantization waste from Theorem~\ref{thm:quantization}: $\sim 31\%$ of binary capacity is wasted encoding continuous partition structures.

A categorical processor simulating physics operates natively in the same partition space as the physical system. There is no discretization error---only the $\eps$-boundary precision, which is a fundamental limit of all bounded computation, not an artifact of binary encoding.

%==============================================================================
\section{Conclusion}
%==============================================================================

We have introduced the categorical converter---a processor that uses virtual membranes to bridge binary digital computation and continuous categorical computation. The key results:

\begin{enumerate}
    \item \textbf{Binary processors are state counters} at base $b = 2$. This is not metaphor---it is the partition-theoretic identity of digital computation.

    \item \textbf{Temperature converts base}: Through partition extinction, thermal control interpolates the effective computational base continuously from binary ($b = 2$) through ternary ($b = 3$) to continuous ($b \to \infty$).

    \item \textbf{Membranes are transducers}: The capacity formula $C(n) = 2n^2$ provides the natural bijection between binary counts and categorical depth. Virtual membranes implement the base conversion physically.

    \item \textbf{The processor is non-quantizable}: The computational unit is $\delta\Depth(T) = 1/\ln b_{\text{eff}}(T)$---a continuous function of temperature with no discrete ``atom'' of computation.

    \item \textbf{Categorical speedup}: S-entropy navigation achieves $O(\log N)$ for problems requiring $O(N)$ binary operations.

    \item \textbf{Biology is categorical}: Cell membranes are natural categorical processors operating at $\sim 10^{20}$ ops/s at the ternary crossover temperature.
\end{enumerate}

The categorical converter unifies digital and analog computation as limiting cases of a single partition-depth framework. It provides the physical foundation for biological computation and suggests new architectures for artificial intelligence. The membrane---far from being a passive barrier---is the most powerful computational element in nature.

\bibliographystyle{unsrt}
\bibliography{references}

\end{document}
